\bibitem{perez2013} Pérez-Rodríguez F, Valero A. Predictive microbiology in foods. New York: Springer; 2013. \url{https://doi.org/10.1007/978-1-4614-5520-2}
\bibitem{tarlak2023} Tarlak F. The use of predictive microbiology for the prediction of the shelf life of food products. Foods. 2023;12:4461. \url{https://doi.org/10.3390/foods12244461}
\bibitem{stavropoulou2019} Stavropoulou E, Bezirtzoglou E. Predictive modeling of microbial behavior in food. Foods. 2019;8:654. \url{https://doi.org/10.3390/foods8120654}
\bibitem{badiamelis2018} Badia-Melis R, Mc Carthy U, Ruiz-Garcia L, Garcia-Hierro J, Robla Villalba JI. New trends in cold chain monitoring applications -- a review. Food Control. 2018;86:170--82. \url{https://doi.org/10.1016/j.foodcont.2017.11.022}
\bibitem{mercier2017} Mercier S, Villeneuve S, Mondor M, Uysal I. Time--temperature management along the food cold chain: a review of recent developments. Compr Rev Food Sci Food Saf. 2017;16:647--67. \url{https://doi.org/10.1111/1541-4337.12269}
\bibitem{jedermann2014} Jedermann R, Nicometo M, Uysal I, et al. Reducing food losses by intelligent food logistics. Philos Trans R Soc A. 2014;372:20130302. \url{https://doi.org/10.1098/rsta.2013.0302}
\bibitem{koutsoumanis2005} Koutsoumanis K, Taoukis PS, Nychas G-JE. Development of a Safety Monitoring and Assurance System for chilled food products. Int J Food Microbiol. 2005;100:253--60. \url{https://doi.org/10.1016/j.ijfoodmicro.2004.10.024}
\bibitem{goransson2018} Göransson M, Nilsson F, Jevinger Å. Temperature performance and food shelf-life accuracy in cold food supply chains: insights from multiple field studies. Food Control. 2018;86:332--41. \url{https://www.sciencedirect.com/science/article/pii/S0956713517305145}
\bibitem{kim2020} Kim J, Santos CA, Kim B-S, Kim J, Koo J. Estimation of real-time remaining shelf life using mean kinetic temperature. LWT. 2020;134:109968. \url{https://doi.org/10.1016/j.lwt.2020.109968}
\bibitem{martin2021} Martin NH, Torres-Frenzel P, Wiedmann M. Invited review: controlling dairy product spoilage to reduce food loss and waste. J Dairy Sci. 2021;104:1251--61. \url{https://doi.org/10.3168/jds.2020-19130}
\bibitem{ranieri2009} Ranieri ML, Boor KJ. Short communication: bacterial ecology of high-temperature, short-time pasteurized milk processed in the United States. J Dairy Sci. 2009;92:4833--40. \url{https://doi.org/10.3168/jds.2009-2181}
\bibitem{who2015} World Health Organization. WHO's first ever global estimates of foodborne diseases find children under 5 account for almost one third of deaths. News release, 3 December 2015. \url{https://www.who.int/news/item/03-12-2015-who-s-first-ever-global-estimates-of-foodborne-diseases-find-children-under-5-account-for-almost-one-third-of-deaths}
\bibitem{fssai2024} Food Safety and Standards Authority of India. Appendix B: microbiological requirements, version 2 (4 November 2024), Tables 2A and 2B. \url{https://fssai.gov.in/upload/uploadfiles/files/Appendix%20B(1).pdf}
\bibitem{lin2016} Lin H, Shavezipur M, Yousef A, Maleky F. Prediction of growth of \textit{Pseudomonas fluorescens} in milk during storage under fluctuating temperature. J Dairy Sci. 2016;99:1822--30. \url{https://doi.org/10.3168/jds.2015-10179}
\bibitem{lima2026} Lima L, Pereira AI, Pintado M, Carocho M. Fresh dairy quality in the era of smart food systems: deterioration pathways, analytical technologies, and predictive modeling. Foods. 2026;15:3439. \url{https://doi.org/10.3390/foods15193439}
\bibitem{gompertz1825} Gompertz B. On the nature of the function expressive of the law of human mortality, and on a new mode of determining the value of life contingencies. Philos Trans R Soc Lond. 1825;115:513--83. \url{https://doi.org/10.1098/rstl.1825.0026}
\bibitem{zwietering1990} Zwietering MH, Jongenburger I, Rombouts FM, van 't Riet K. Modeling of the bacterial growth curve. Appl Environ Microbiol. 1990;56:1875--81. \url{https://doi.org/10.1128/aem.56.6.1875-1881.1990}
\bibitem{baranyi1993} Baranyi J, Roberts TA, McClure P. A non-autonomous differential equation to model bacterial growth. Food Microbiol. 1993;10:43--59. \url{https://doi.org/10.1006/fmic.1993.1005}
\bibitem{baranyi1994} Baranyi J, Roberts TA. A dynamic approach to predicting bacterial growth in food. Int J Food Microbiol. 1994;23:277--94. \url{https://doi.org/10.1016/0168-1605(94)90157-0}
\bibitem{ratkowsky1982} Ratkowsky DA, Olley J, McMeekin TA, Ball A. Relationship between temperature and growth rate of bacterial cultures. J Bacteriol. 1982;149:1--5. \url{https://doi.org/10.1128/jb.149.1.1-5.1982}
\bibitem{rosso1993} Rosso L, Lobry JR, Flandrois JP. An unexpected correlation between cardinal temperatures of microbial growth highlighted by a new model. J Theor Biol. 1993;162:447--63. \url{https://doi.org/10.1006/jtbi.1993.1099}
\bibitem{taoukis1989} Taoukis PS, Labuza TP. Applicability of time-temperature indicators as shelf life monitors of food products. J Food Sci. 1989;54:783--8. \url{https://doi.org/10.1111/j.1365-2621.1989.tb07882.x}
\bibitem{baranyi2004} Baranyi J, Tamplin ML. ComBase: a common database on microbial responses to food environments. J Food Prot. 2004;67:1967--71. \url{https://doi.org/10.4315/0362-028X-67.9.1967}
\bibitem{combase} ComBase. Database of microbial responses to food environments: growth records with milk as the food matrix. Accessed 19 Oct 2026. \url{https://www.combase.cc}
\bibitem{neumeyer1995} Neumeyer K. Modelling pseudomonad growth in milk and milk-based products. Research Master thesis. Hobart: University of Tasmania; 1995. \url{https://eprints.utas.edu.au/21013/}
\bibitem{neumeyer1997a} Neumeyer K, Ross T, McMeekin TA. Development of a predictive model to describe the effects of temperature and water activity on the growth of spoilage pseudomonads. Int J Food Microbiol. 1997;38:45--54. \url{https://doi.org/10.1016/S0168-1605(97)00089-5}
\bibitem{corradini2005} Corradini MG, Peleg M. Estimating non-isothermal bacterial growth in foods from isothermal experimental data. J Appl Microbiol. 2005;99:187--200. \url{https://doi.org/10.1111/j.1365-2672.2005.02570.x}
\bibitem{zwietering1994} Zwietering MH, de Wit JC, Cuppers HGAM, van 't Riet K. Modeling of bacterial growth with shifts in temperature. Appl Environ Microbiol. 1994;60:204--13. \url{https://doi.org/10.1128/aem.60.1.204-213.1994}
\bibitem{koutsoumanis2001} Koutsoumanis K. Predictive modeling of the shelf life of fish under nonisothermal conditions. Appl Environ Microbiol. 2001;67:1821--9. \url{https://doi.org/10.1128/AEM.67.4.1821-1829.2001}
\bibitem{ndraha2018} Ndraha N, Hsiao H-I, Vlajic J, Yang M-F, Lin H-TV. Time-temperature abuse in the food cold chain: review of issues, challenges, and recommendations. Food Control. 2018;89:12--21. \url{https://doi.org/10.1016/j.foodcont.2018.01.027}
\bibitem{zou2023} Zou Y, Wu J, Wang X, Morales K, Liu G, Manzardo A. An improved artificial neural network using multi-source data to estimate food temperature during multi-temperature delivery. J Food Eng. 2023;351:111518. \url{https://doi.org/10.1016/j.jfoodeng.2023.111518}
\bibitem{muller2026} Müller WA, Kronbauer M, Ferreira SB, da Silva SB. Computational assessment of temperature failures in the cold supply chain: implications for beef shelf life. Food Control. 2026;187:112100. \url{https://doi.org/10.1016/j.foodcont.2026.112100}
\bibitem{sofat2026} Sofat A, Sodhi B. Beyond thresholds: a quality-aware decision intelligence framework for cold chain IoT systems. arXiv preprint arXiv:2608.15082. 2026. \url{https://arxiv.org/abs/2608.15082}
\bibitem{nauta2000} Nauta MJ. Separation of uncertainty and variability in quantitative microbial risk assessment models. Int J Food Microbiol. 2000;57:9--18. \url{https://doi.org/10.1016/S0168-1605(00)00225-7}
\bibitem{lau2022} Lau S, Trmčić A, Martin NH, Wiedmann M, Murphy SI. Development of a Monte Carlo simulation model to predict pasteurized fluid milk spoilage due to post-pasteurization contamination with gram-negative bacteria. J Dairy Sci. 2022;105:1978--98. \url{https://doi.org/10.3168/jds.2021-21316}
\bibitem{buehler2018} Buehler AJ, Martin NH, Boor KJ, Wiedmann M. Psychrotolerant spore-former growth characterization for the development of a dairy spoilage predictive model. J Dairy Sci. 2018;101:6964--81. \url{https://doi.org/10.3168/jds.2018-14501}
\bibitem{qian2023} Qian C, Murphy SI, Lott TT, Martin NH, Wiedmann M. Development and deployment of a supply-chain digital tool to predict fluid-milk spoilage due to psychrotolerant sporeformers. J Dairy Sci. 2023;106:8414--33. \url{https://doi.org/10.3168/jds.2023-23673}
\bibitem{giannakourou2020} Giannakourou MC, Taoukis PS. Holistic approach to the uncertainty in shelf life prediction of frozen foods at dynamic cold chain conditions. Foods. 2020;9:714. \url{https://doi.org/10.3390/foods9060714}
\bibitem{ds18b20} Analog Devices (Maxim Integrated). DS18B20 programmable resolution 1-Wire digital thermometer. Product page and datasheet. \url{https://www.analog.com/en/DS18B20}
\bibitem{iec60751} International Electrotechnical Commission. IEC 60751:2022 Industrial platinum resistance thermometers and platinum temperature sensors. Geneva: IEC; 2022. \url{https://webstore.iec.ch/en/publication/63753}
\bibitem{hanley1982} Hanley JA, McNeil BJ. The meaning and use of the area under a receiver operating characteristic (ROC) curve. Radiology. 1982;143:29--36. \url{https://doi.org/10.1148/radiology.143.1.7063747}
\bibitem{mannwhitney1947} Mann HB, Whitney DR. On a test of whether one of two random variables is stochastically larger than the other. Ann Math Stat. 1947;18:50--60. \url{https://doi.org/10.1214/aoms/1177730491}
\bibitem{efron1993} Efron B, Tibshirani RJ. An introduction to the bootstrap. New York: Chapman \& Hall; 1993. \url{https://doi.org/10.1201/9780429246593}
\bibitem{delong1988} DeLong ER, DeLong DM, Clarke-Pearson DL. Comparing the areas under two or more correlated receiver operating characteristic curves: a nonparametric approach. Biometrics. 1988;44:837--45. \url{https://doi.org/10.2307/2531595}
\bibitem{brier1950} Brier GW. Verification of forecasts expressed in terms of probability. Mon Weather Rev. 1950;78:1--3. \url{https://doi.org/10.1175/1520-0493(1950)078<0001:VOFEIT>2.0.CO;2}
\bibitem{naeini2015} Naeini MP, Cooper GF, Hauskrecht M. Obtaining well calibrated probabilities using Bayesian binning. Proc AAAI Conf Artif Intell. 2015;29:2901--7. \url{https://doi.org/10.1609/aaai.v29i1.9602}
\bibitem{ross1996} Ross T. Indices for performance evaluation of predictive models in food microbiology. J Appl Bacteriol. 1996;81:501--8. \url{https://doi.org/10.1111/j.1365-2672.1996.tb03539.x}
\bibitem{harris2020} Harris CR, Millman KJ, van der Walt SJ, et al. Array programming with NumPy. Nature. 2020;585:357--62. \url{https://doi.org/10.1038/s41586-020-2649-2}
\bibitem{virtanen2020} Virtanen P, Gommers R, Oliphant TE, et al. SciPy 1.0: fundamental algorithms for scientific computing in Python. Nat Methods. 2020;17:261--72. \url{https://doi.org/10.1038/s41592-019-0686-2}
\bibitem{hunter2007} Hunter JD. Matplotlib: a 2D graphics environment. Comput Sci Eng. 2007;9:90--5. \url{https://doi.org/10.1109/MCSE.2007.55}
